\section{Theory}
\label{sec:theory}

The propositions of this section state why the protocol should work and under which conditions; each is checked numerically in Section~\ref{sec:results}, and the conditions under which it fails are stated with it. Proposition~\ref{prop:unbiased} says what the holdout estimate is when no step read the holdout. Propositions~\ref{prop:detection} and~\ref{prop:discrete} say with what probability the third gate detects a leak through a quantile edge and through a discrete choice. Proposition~\ref{prop:selection} bounds the optimism of a choice made on the holdout. Propositions~\ref{prop:leakauc} and~\ref{prop:psi} describe a leak that exists only in development and explain what the stability index sees of it. The scaling of the cross-validation bias is stated as a conjecture.

\subsection{The holdout estimate when no step read the holdout}

\begin{proposition}[Unbiasedness]
\label{prop:unbiased}
Let $\theta=\mathcal A(D_R)$ with $R\cap H=\emptyset$, and let the rows of $H$ be an independent sample from the population, independent of $D_R$, with $n_1\ge1$ events and $n_0\ge1$ non-events. Then, conditionally on $\theta$ and on the labels of $H$, the holdout statistic of equation~\eqref{eq:auchat} satisfies $\E[\hat A_H\mid\theta,y_H]=A^{\ast}(\theta)$. If $s^2$ is the DeLong estimator of the variance of $\hat A_H$, then $(\hat A_H-A^{\ast}(\theta))/s$ converges in distribution to $\mathcal N(0,1)$ as $n_1,n_0\to\infty$, and the interval $\hat A_H\pm1.96\,s$ has asymptotic coverage $0.95$.
\end{proposition}
\begin{proof}
Given $\theta$, the score function $S_\theta$ is fixed. Given the labels of $H$, the inputs of the events are independent draws from the distribution of $X$ given $y=1$ and those of the non-events from that given $y=0$, independent of each other. For one event $i$ and one non-event $j$ the kernel $\psi_{ij}=\indic\{S_\theta(x_i)>S_\theta(x_j)\}+\frac12\indic\{S_\theta(x_i)=S_\theta(x_j)\}$ has expectation $A^{\ast}(\theta)$ by equation~\eqref{eq:aucstar}. The statistic $\hat A_H$ is the average of $\psi_{ij}$ over the $n_1n_0$ pairs, so its conditional expectation is $A^{\ast}(\theta)$. It is a two-sample $U$-statistic with a bounded kernel, whose asymptotic normality and whose consistent variance estimator by structural components are the results behind the DeLong interval \citep{delong1988,hanley1982}.
\end{proof}
The proposition fails exactly when $R\cap H\neq\emptyset$: then $\theta$ is a function of the rows of $H$, the pairs are no longer independent of $S_\theta$, and the conditional expectation of the statistic is no longer $A^{\ast}(\theta)$. The sign of the resulting bias is not determined by the proposition; for a procedure that adapts to the rows it reads, it is positive, and Section~\ref{sec:results} measures its size for nine ways of reading the holdout. The proposition also fails when assumption A1 fails: a holdout from a later vintage estimates the performance on that vintage, not on the population of the training rows.

\subsection{Detection by the outside-information gate}

\begin{proposition}[Detection of a leak through a quantile edge]
\label{prop:detection}
Let a step fit the quantile edge at level $q$ of $N$ pooled, continuous, exchangeable values, with $(N-1)q$ not an integer, by linear interpolation between the order statistics of ranks $k-1$ and $k$ with $k=\lfloor(N-1)q\rfloor+2$, and let this edge be read from a pool that contains $m$ outside rows. Let the third gate send these $m$ rows above the maximum of the pool, the $m$ rows being a uniformly random subset of the $N$ pooled rows. Then the edge changes unless every moved row has rank above $k$ in the pool, and
\begin{equation}
\begin{split}
\Prob(\text{gate detects})&=1-\frac{\binom{N-k}{m}}{\binom{N}{m}}\\
&\ge\ 1-(1-q)^m .
\end{split}
\label{eq:detect}
\end{equation}
\end{proposition}
\begin{proof}
The order statistic of rank $j\le N-m$ of the perturbed pool is the $j$-th smallest of the $N-m$ rows that were not moved. It equals the original order statistic of rank $j$ if and only if no moved row has original rank at most $j$. The edge is a function of the order statistics of ranks $k-1$ and $k$, so it is unchanged if no moved row has rank at most $k$; if a moved row has rank $r\le k$, every order statistic of rank $j\ge r$ is replaced by a strictly larger value, in particular the one of rank $k$, and its weight in the interpolation, $(N-1)q-\lfloor(N-1)q\rfloor$, is positive, so the edge changes. The moved set is a uniformly random $m$-subset of the ranks, so the probability that all its elements lie among the $N-k$ ranks above $k$ is $\binom{N-k}{m}/\binom{N}{m}=\prod_{j=0}^{m-1}\frac{N-k-j}{N-j}$. Each factor is at most $\frac{N-k}{N}$, and $k\ge(N-1)q+1$ gives $\frac{N-k}{N}\le\frac{(N-1)(1-q)}{N}<1-q$.
\end{proof}
For the pool of the experiments, $N=\nv{th.det.ntrain}+\nv{th.det.nout}$ pools the development rows (the training and the calibration slices) and the holdout rows and $q=0.9$, so that $k=1351$, the probability that one moved row is detected is $\nv{th.det.1.exact}\%$ and that two are detected is $\nv{th.det.2.exact}\%$. The gate that moves the whole holdout, $m=|H|=450$, cannot miss, since $m$ exceeds $N-k=149$ and the probability of the proposition is exactly zero, so that the detection of a leak through a step that fits an edge on pooled rows does not depend on the gate but on whether the step reads the rows at all. A step that fits a median or a mean has the same property with $q=\frac12$ or, for a mean, with certainty.

\subsection{Optimism of a choice made on the holdout}

\begin{proposition}[Selection on the holdout]
\label{prop:selection}
Let $K$ configurations have true values $a_k=A^{\ast}(\theta_k)$ and holdout estimates $\hat A_k=a_k+Z_k$, where each $Z_k$ is sub-Gaussian with variance proxy $\sigma^2$, that is $\E e^{\lambda Z_k}\le e^{\lambda^2\sigma^2/2}$. Let $\hat k=\arg\max_k\hat A_k$. Then
\begin{equation}
\E\bigl[\hat A_{\hat k}-a_{\hat k}\bigr]\ \le\ \sigma\sqrt{2\ln K}.
\label{eq:selectbound}
\end{equation}
For the Mann-Whitney statistic of equation~\eqref{eq:auchat} one may take $\sigma^2=1/(4\min(n_1,n_0))$.
\end{proposition}
\begin{proof}
$\hat A_{\hat k}-a_{\hat k}=Z_{\hat k}\le\max_kZ_k$. For $\lambda>0$, Jensen's inequality gives $\E\max_kZ_k\le\frac1\lambda\ln\E\sum_ke^{\lambda Z_k}\le\frac{\ln K}{\lambda}+\frac{\lambda\sigma^2}{2}$, and the choice $\lambda=\sqrt{2\ln K}/\sigma$ gives the bound. The statistic $\hat A_H$ is a two-sample $U$-statistic with a kernel in $[0,1]$; Hoeffding's representation of it as an average of means of $\min(n_1,n_0)$ independent pairs and the convexity of the exponential give $\E\exp\{\lambda(\hat A_H-\E\hat A_H)\}\le\exp\{\lambda^2/(8\min(n_1,n_0))\}$, which is the condition on $Z_k$ with $\sigma^2=1/(4\min(n_1,n_0))$ \citep{hoeffding1963,boucheron2013}.
\end{proof}
The bound is for the worst case, in which the estimates are independent. Configurations of one pipeline are strongly dependent, and the maximum of $K$ standard normal variables with pairwise correlation $\rho$ has a mean of $\nv{th.max.8.r0.0}$ for $K=8$ at $\rho=0$, $\nv{th.max.8.r0.5}$ at $\rho=0.5$ and $\nv{th.max.8.r0.9}$ at $\rho=0.9$, in units of $\sigma$, against the bound $\nv{th.max.8.bound}$ (simulation with 20,000 draws per cell, the correlation induced by a common factor). The optimism of a choice made on the holdout is therefore small when the configurations are similar and grows with $\sqrt{\ln K}$ when they are not, which is the pattern of Section~\ref{sec:results}.

\subsection{A leak that exists only in development}

The planted leak of the engine's generator is the case in which a feature carries the outcome in development and is noise at decision time. Let the score be $S=w_1G+w_2F$, where $G\sim\mathcal N(d\,y,1)$ is the genuine feature and $F$ an independent feature that is $\mathcal N(s\,y,1)$ in development and $\mathcal N(0,1)$ at decision time.

\begin{proposition}[AUC of a score with a planted leak]
\label{prop:leakauc}
In development, and at decision time, the AUC of $S$ is
\begin{equation}
\AUC_{\mathrm{dev}}=\Phi\!\left(\frac{w_1d+w_2s}{\sqrt{2(w_1^2+w_2^2)}}\right),\qquad
\AUC_{\mathrm{cur}}=\Phi\!\left(\frac{w_1d}{\sqrt{2(w_1^2+w_2^2)}}\right),
\label{eq:leakauc}
\end{equation}
where $\Phi$ is the standard normal distribution function. The development AUC is largest for $w\propto(d,s)$, where it equals $\Phi\bigl(\sqrt{(d^2+s^2)/2}\bigr)$, and the decision-time AUC of that score is $\Phi\bigl(d^2/\sqrt{2(d^2+s^2)}\bigr)$, which is below the AUC $\Phi(d/\sqrt2)$ of the genuine feature alone whenever $s>0$. For equal weights, $\AUC_{\mathrm{dev}}=\Phi((d+s)/2)$ and $\AUC_{\mathrm{cur}}=\Phi(d/2)$.
\end{proposition}
\begin{proof}
The score difference $S_1-S_0$ between a defaulter and a survivor is normal with mean $w_1d+w_2s$ in development, $w_1d$ at decision time, and variance $2(w_1^2+w_2^2)$, so the AUC $\Prob(S_1-S_0>0)$ is $\Phi(\mu/\sqrt{2(w_1^2+w_2^2)})$ with the mean $\mu$ of the case. By the Cauchy-Schwarz inequality $w_1d+w_2s\le\sqrt{d^2+s^2}\sqrt{w_1^2+w_2^2}$, with equality at $w\propto(d,s)$, which gives the maximum; substituting $w_1/\sqrt{w_1^2+w_2^2}=d/\sqrt{d^2+s^2}$ gives the decision-time AUC. The AUC of the genuine feature alone is the case $w_2=0$. The comparison follows from $d/\sqrt{d^2+s^2}<1$.
\end{proof}
At $d=1$ and $s=1$ with equal weights, the development AUC is $\Phi(1)=\nv{e3.s1.0.devcf}$ and the decision-time AUC $\Phi(\tfrac12)=\nv{e3.s1.0.curcf}$: the leak inflates the development AUC by $\nv{e3.s1.0.infl}$, and the model is worse at decision time than the genuine feature alone, whose AUC is $\Phi(1/\sqrt2)=\nv{th.opt.1.0.clean}$. The test that compares the current AUC with its development value sees the whole gap. The index that compares the distribution of the leaked feature between development and a reference sees a much smaller quantity.

\begin{proposition}[Stability index of the leaked feature]
\label{prop:psi}
Let the leaked feature be $\mathcal N(0,1)$ for survivors and $\mathcal N(s,1)$ for the share $\pi$ of defaulters, and let the reference sample be $\mathcal N(0,1)$. For a binning with reference probabilities $p_b$ and defaulter probabilities $r_b=\Phi(e_b-s)-\Phi(e_{b-1}-s)$, the population stability index of the mixture is
\begin{equation}
\mathrm{PSI}=\pi^2\sum_b\frac{(r_b-p_b)^2}{p_b}+O(\pi^3),\qquad\text{and, for fine bins, }\ \mathrm{PSI}\approx\pi^2\bigl(e^{s^2}-1\bigr).
\label{eq:psi}
\end{equation}
\end{proposition}
\begin{proof}
The mixture has bin probabilities $m_b=p_b+\pi(r_b-p_b)$. With $x_b=\pi(r_b-p_b)/p_b$ the index is $\sum_b(m_b-p_b)\ln(m_b/p_b)=\sum_bp_bx_b\ln(1+x_b)=\sum_bp_b\bigl(x_b^2-\tfrac12x_b^3+\dots\bigr)$, and the quadratic term is the first term of~\eqref{eq:psi}. For fine bins $\sum_b(r_b-p_b)^2/p_b$ tends to the chi-square divergence between $\mathcal N(s,1)$ and $\mathcal N(0,1)$, which is $e^{s^2}-1$.
\end{proof}
The index of two samples of 5,000 rows each, binned in ten deciles, has an expectation of $(B-1)(1/5000+1/5000)=\nv{th.psi.null}$ when the two samples agree \citep{yurdakul2020}. With the generator's share $\pi=0.03$ the second-order form gives $\nv{th.psi.1.0.binned_chi2_form}$ at $s=1$, below that level: the leak is invisible to the index. The binned form is within 11\% of the exact index of the engine at every $s$ up to 3 (\nv{th.psi.3.0.binned_chi2_form} against \nv{th.psi.3.0.engine_psi} at $s=3$); the continuous form is within 13\% up to $s=0.5$ and diverges beyond it (\nv{th.psi.3.0.continuous_form} at $s=3$), because the decile grid cannot resolve the tails that the exponential describes. The simulated index under no leak has a mean of \nv{e3.s0.0.psi}, which agrees with the expectation.

\subsection{Detection of a discrete choice by repeated perturbation}

The gates of the protocol perturb the holdout once by default, and a leak that consists in choosing among a few configurations on the holdout leaves a discrete trace in the fitted parameters: the parameters change only when the perturbed holdout makes another configuration the best one. The next proposition gives the detection probability of that trace as a function of the number of perturbation draws.

\begin{proposition}[Detection of a discrete choice]
\label{prop:discrete}
Let a leak consist in choosing one of $K$ configurations by the highest score on $H$, and let the perturbation of $H$ be such that the choice that would be made on the perturbed holdout is uniformly distributed on the $K$ configurations, independently across $d$ independent perturbation draws. Then the third gate fails, for the base choice, with probability $1-K^{-d}$.
\end{proposition}
\begin{proof}
A draw does not change the choice with probability $1/K$, because the perturbed choice is uniform and equals the base choice with that probability. The draws are independent, so the choice survives all $d$ draws with probability $K^{-d}$.
\end{proof}
Uniformity is the condition. When the labels of $H$ are redrawn at random the scores of the configurations on $H$ are exchangeable if the configurations are similar models, and the proposition applies; when one configuration dominates for structural reasons, the perturbed choice equals the base choice more often than $1/K$ and the detection is lower than the proposition states. The measurement of Section~\ref{sec:results} is the check.

\subsection{Scaling of the bias of cross-validation with preprocessing outside the folds}

\begin{conjecture}[Scaling of the cross-validation bias]
\label{conj:cv}
For a scorecard whose binning, weight of evidence and screening of the best $k$ of $p$ inputs are fitted on all development rows before the folds, the bias of the five-fold AUC against the AUC of the fold models has the form $b(n,p)\approx c(p)\,n^{-\alpha}$ with $\alpha$ between $0.8$ and $1$, and $c(p)$ increasing and concave in $p$.
\end{conjecture}
The heuristic is that the optimism of an in-sample AUC grows with the effective number of fitted parameters divided by $n$, which gives $\alpha=1$, and that the screening of the best ten inputs among $p$ adds an optimism that grows with $p$ but saturates once the ten best noise inputs are essentially the extreme order statistics. The evidence is Table~\ref{tab:cv-scaling}: a log-log regression over $n=400$ to $6{,}400$ at $p=40$ gives a slope of $-\nv{e5.exp}\pm\nv{e5.expse}$ (standard error), and at $n=1{,}600$ the bias moves from \nv{e5.n1600.p10.out} at 10 inputs to \nv{e5.n1600.p160.out} at 160. The slope is below one in magnitude by more than ten standard errors, so the conjecture is stated with $\alpha<1$; the concavity in $p$ is read from five points and is weak evidence. The range of $\alpha$ between 0.8 and 1 was read from this same fit, so that the experiment describes the conjecture and does not test it; a test needs data sets or cells that were not used to state it. The conjecture is not proved, and it is stated for the scorecard of the experiments, whose informative inputs are five of $p$.
